\documentclass{article}
\usepackage[utf8]{inputenc}
\usepackage{amsmath,amssymb}
\usepackage{booktabs}
\usepackage{graphicx}
\usepackage{hyperref}
\usepackage[margin=1in]{geometry}
\usepackage{algorithm}
\usepackage{algpseudocode}

\title{FoveatedKV: Fused Importance-Adaptive KV Cache Compression\\with Spike-Driven Promotion for Apple Silicon}
\author{Sam Furr}
\date{}

\begin{document}
\maketitle

\begin{abstract}
We present FoveatedKV, a KV cache compression system for LLM inference
on Apple Silicon that combines three ideas: (1)~importance-adaptive
two-tier quantization inspired by foveated rendering, keeping the top
10\% of tokens at fp16 and compressing the rest with fp8~E4M3 keys and
INT4 values; (2)~a fused Split-K Metal compute kernel that reads
quantized data directly from unified memory, dequantizes in registers,
and performs spike detection as a free byproduct of online softmax; and
(3)~a C++ promotion pipeline that recovers exact fp16 representations
from an NVMe-backed archive when the kernel detects a far-tier token
receiving unexpectedly high attention. The compressed cache blob is tracked
as a proper MLX input array for correct graph pipelining. The asymmetric
K/V precision strategy (higher precision for keys than values) builds on
prior work by KIVI, LeanKV, and others. Our contribution is a complete
system that implements this strategy within an importance-adaptive
framework, targeting the memory-bandwidth constraints of Apple Silicon.
On an 8GB Mac, FoveatedKV achieves 2.02$\times$ memory compression with
0.995+ cosine fidelity. At the kernel level, the fused kernel reaches
3.34$\times$ speedup over Apple's SDPA at 16K context on 7B-equivalent
shapes. On memory-constrained hardware (7B model on 8GB), the compressed
cache makes foveated 2.3$\times$ faster than standard inference by
reducing memory pressure. The system recovers facts lost to quantization
via spike-driven promotion, matching fp16 retrieval accuracy.
\end{abstract}

\section{Introduction}

Large language model inference at long context is bottlenecked by memory
bandwidth, not compute. Every decode step reads the entire KV cache from
memory. For a 7B model at 32K context, the KV cache alone exceeds 1\,GB.
On consumer hardware with limited memory bandwidth (e.g., Apple M-series
with 100--200\,GB/s), this dominates inference latency.

Existing approaches to KV cache compression fall into two categories:
\emph{eviction} methods (H2O~\cite{h2o}, SnapKV~\cite{snapkv},
StreamingLLM~\cite{streamingllm}) that discard low-importance tokens,
and \emph{quantization} methods (KIVI~\cite{kivi},
LeanKV~\cite{leankv}, KV-AdaQuant~\cite{kvadaquant}) that reduce
precision. Eviction methods break the softmax denominator by removing
tokens from the sum. Uniform quantization treats high-attention tokens
identically to filler text.

Recent work has established that keys and values have different
quantization sensitivity: key errors are amplified through $\exp()$ in
softmax, while value errors are bounded by attention weight. This
observation has been independently reported by
KIVI~\cite{kivi}, LeanKV~\cite{leankv}, KV-AdaQuant~\cite{kvadaquant},
and AsymKV~\cite{asymkv}, leading to proposals for asymmetric precision
(e.g., K8V4). We build on this established insight.

Our contribution is \emph{systems-level}: a complete implementation that
combines asymmetric K/V quantization with importance-adaptive tiering,
a fused GPU kernel, and a spike-driven promotion mechanism---all
targeting the specific constraints of Apple Silicon. Concretely:

\begin{enumerate}
\item \textbf{Foveated two-tier cache.} Tokens are partitioned per
  attention head into a near tier (fp16) and far tier (fp8 K + INT4 V)
  based on attention importance. Every token remains in the softmax
  computation---no eviction, no broken denominators.
\item \textbf{Fused Split-K Metal kernel.} A single GPU dispatch reads
  mixed-precision data from unified memory, dequantizes in registers
  (fp8 via 256-entry LUT, INT4 via shift-and-mask), computes online
  softmax, and detects spikes---all without materializing fp16
  intermediates to global memory.
\item \textbf{Spike detection and promotion.} The kernel identifies
  far-tier tokens receiving unexpectedly high attention scores as a free
  byproduct of online softmax tracking. A C++ background worker reads
  exact fp16 originals from an NVMe-backed mmap archive and promotes
  them into the near tier, recovering quality lost to quantization.
\end{enumerate}

\section{Background: Asymmetric K/V Precision}
\label{sec:asymmetric}

Multiple recent works have observed that keys and values have
fundamentally different error sensitivity in the attention computation.
We briefly review the theoretical basis, as it motivates our design
choices.

Consider the attention output for a single query $q$:
\begin{equation}
  \text{out} = \sum_i \frac{\exp(q \cdot k_i / \sqrt{d})}{\sum_j \exp(q \cdot k_j / \sqrt{d})} \cdot v_i
  = \sum_i \alpha_i v_i
\end{equation}

\paragraph{Key error amplification.}
A quantization error $\epsilon_k$ on key $k_i$ produces a score error
$\delta_s = q \cdot \epsilon_k / \sqrt{d}$. Through $\exp()$, this
becomes a multiplicative factor on the attention weight:
\begin{equation}
  \alpha_i' = \alpha_i \cdot \exp(\delta_s)
\end{equation}
Because softmax normalizes, this error redistributes weight across
\emph{all} tokens in the sequence. KIVI~\cite{kivi} identified this
through the lens of per-channel vs per-token quantization;
LeanKV~\cite{leankv} proposed explicit K8V4 configurations;
KV-AdaQuant~\cite{kvadaquant} provided systematic analysis via singular
value distributions and spectral norms; AsymKV~\cite{asymkv} conducted
a systematic examination of attention output error.

\paragraph{Value error is bounded.}
A quantization error $\epsilon_v$ on value $v_i$ contributes
$\alpha_i \cdot \epsilon_v$ to the output. For low-attention tokens,
$\alpha_i$ is small, so the output error is bounded:
\begin{equation}
  \|\text{error}\| \leq \sum_{i \in \text{far}} \alpha_i \|\epsilon_{v,i}\| \leq \alpha_{\text{far}} \cdot \max_i \|\epsilon_{v,i}\|
\end{equation}

\paragraph{Our choice: fp8 E4M3 keys + INT4 values.}
Following LeanKV's K8V4 configuration, we use fp8~E4M3 for far-tier
keys and INT4 for far-tier values. We prefer fp8 over INT8 for keys
because the 256-entry lookup table (all fp8$\to$fp16 mappings
precomputed at kernel dispatch time) enables a single shared-memory read
per element, replacing $\sim$10 ALU ops for per-channel INT8
dequantization. Our ablation (Table~\ref{tab:key_ablation}) confirms
that fp8 and INT8 achieve equivalent error, while INT4 keys produce
$34\times$ more error.

We empirically validate the asymmetric choice on our system:

\begin{table}[h]
\centering
\begin{tabular}{lcc}
\toprule
Configuration & Cosine vs Exact & Cosine Error \\
\midrule
fp8 K + INT4 V (ours) & 0.999939 & 0.000061 \\
INT4 K + INT4 V (symmetric) & 0.999783 & 0.000217 \\
\bottomrule
\end{tabular}
\caption{Asymmetric vs symmetric K/V precision at 4K context.
Symmetric INT4 produces $3.6\times$ more attention error, consistent
with findings in prior work~\cite{leankv, kvadaquant}.}
\label{tab:asymmetric}
\end{table}

\section{System Design}

While the quantization strategy builds on prior work, the system
architecture is our primary contribution: a foveated cache with a fused
kernel and promotion pipeline, designed for Apple Silicon's unified
memory.

\subsection{Two-Tier Foveated Cache}

Inspired by foveated rendering in VR~\cite{guenter2012}---where only the
gaze center is rendered at full resolution---we partition tokens into two
tiers per attention head after prefill:

\begin{itemize}
\item \textbf{Near tier} (default 10\%): attention sinks (first $N$
  tokens) + recency window + top importance tokens, stored at fp16.
\item \textbf{Far tier} (90\%): remaining tokens, stored with fp8~E4M3
  keys and INT4 (4-bit asymmetric per-token) values.
\end{itemize}

The ``gaze'' is the query vector; the ``fovea'' is the near tier; the
``periphery'' is the far tier. Unlike eviction methods, every token
remains in the attention computation. The softmax denominator is exact.
The only approximation is quantization noise on low-attention tokens.

This importance-adaptive approach differs from uniform quantization
methods like KIVI~\cite{kivi}: high-attention tokens retain full
precision, concentrating quantization noise where it has minimal impact.
LeanKV~\cite{leankv} also combines importance-adaptive compression with
asymmetric precision, but targets CUDA GPUs with a different kernel
design. Our system targets Apple Silicon's unified memory architecture,
which enables zero-copy access patterns that discrete-GPU systems cannot
exploit.

\subsection{Fused Split-K Metal Kernel}

We implement a fused Metal compute kernel that performs the entire
attention computation---across both tiers and the decode buffer---in a
single GPU dispatch:

\begin{enumerate}
\item Load query, pre-scale by $1/\sqrt{d}$.
\item Near tier: standard fp16 dot product + online softmax accumulation.
\item Far tier: fp8 key dequantization via 256-entry threadgroup LUT
  (all 256 fp8$\to$fp16 mappings precomputed at dispatch time---one
  shared-memory read replaces $\sim$10 ALU ops per element),
  dot product, score-gated INT4 value dequantization (skip when
  $\text{score} < m - 16$, i.e., $\exp(-16) \approx 10^{-7}$).
\item Decode buffer: new tokens added since compression (fp16).
\item Shared-memory reduction across Split-K groups.
\item Spike detection: track $\max(\text{far scores})$ vs
  $\min(\text{near scores})$ per head as a free byproduct of online
  softmax.
\end{enumerate}

The kernel never materializes fp16 intermediates to global memory.
Dequantization happens in registers. This is where the bandwidth
savings come from: reading fp8+INT4 data costs 2--4$\times$ fewer bytes
than fp16.

Key kernel optimizations: (a)~pre-scaled query amortizes the
$1/\sqrt{d}$ multiply across the entire sequence; (b)~single-exp online
softmax computes each $\exp()$ exactly once with FMA for rescale +
accumulate; (c)~4-token blocked far loop provides instruction-level
parallelism across independent K dot products; (d)~score-gated V
loading skips INT4 dequantization for negligible-attention tokens.

\paragraph{TurboQuant (experimental).} An optional mode doubles the
compression ratio to 4$\times$ by using INT4 for both keys and values
in the far tier. This sacrifices key precision (see
Table~\ref{tab:key_ablation}) for additional bandwidth savings. We
include it as an experimental option for scenarios where throughput is
prioritized over fidelity.

\subsection{Spike Detection and Promotion}
\label{sec:promotion}

The spike detection mechanism is a distinguishing feature of our system.
While LeanKV~\cite{leankv} and KIVI~\cite{kivi} apply static
quantization, FoveatedKV can dynamically recover precision when needed.

During online softmax, the kernel naturally tracks $\max(\text{far
scores})$ and $\min(\text{near scores})$ per head. When a far-tier token
scores above the near-tier minimum, it is flagged as a \emph{spike}.
This detection adds zero computational overhead---the max/min values are
already maintained by the softmax accumulator.

A C++ \texttt{PromotionPipeline} background worker handles promotion:

\begin{enumerate}
\item Reads the exact fp16 original from an NVMe-backed mmap archive
  (the archive preserves all pre-quantization KV pairs on disk).
\item Writes it into pre-allocated headroom slots in the near tier blob.
\item Atomically increments the per-head valid count.
\end{enumerate}

The kernel sees the promoted token on the next dispatch as an ordinary
near-tier token. Promotion is lossless---the disk archive preserves
exact fp16 originals---and the unified memory architecture of Apple
Silicon enables zero-copy access between the CPU promotion worker and
the GPU kernel.

\paragraph{Filtering is critical.}
Raw spike rates are high (95\% of head-layer slots per step), so the
pipeline applies aggressive filtering: per-(layer,head) cooldown,
splitmix64 position dedup, per-step budget cap, and GQA head dedup.
Of 15,748 raw spikes over 80 decode steps, only 494 (3.1\%) result in
actual promotions. Without filtering, every far token would be promoted,
defeating tiered compression.

\section{Results}

All experiments run on Apple M2 8GB. We evaluate on two model families:
Qwen2.5-0.5B and Llama-3.2-1B. Kernel microbenchmarks use 7B-equivalent
tensor shapes (H$_\text{kv}$=8, H$_\text{q}$=32, D=128).

\subsection{Kernel Performance}

Fused kernel vs Apple's SDPA using 7B-equivalent tensor shapes
(H$_\text{kv}$=8, H$_\text{q}$=32, D=128; these are the GQA
dimensions of models like Qwen2.5-7B and Llama-3-8B), single layer:

\begin{table}[h]
\centering
\begin{tabular}{rccc}
\toprule
Context & fp16 SDPA & Fused Kernel & Speedup \\
\midrule
1K & 0.84\,ms & 1.00\,ms & 0.84$\times$ \\
4K & 2.07\,ms & 1.20\,ms & 1.72$\times$ \\
8K & 4.15\,ms & 1.68\,ms & 2.47$\times$ \\
16K & 9.67\,ms & 2.90\,ms & 3.34$\times$ \\
32K & 15.19\,ms & 5.18\,ms & 2.93$\times$ \\
\bottomrule
\end{tabular}
\caption{Fused kernel speedup over Apple SDPA. The fused kernel has
higher fixed overhead, so SDPA is faster at 1K. Above 4K, the
bandwidth advantage dominates and grows with context length, reaching
3.34$\times$ at 16K.}
\end{table}

\subsection{End-to-End Throughput}

On the small 0.5B model, foveated decode is somewhat slower than standard
due to Python SDPA interceptor overhead. The kernel itself is faster
(Table~\ref{tab:kernel}), but per-layer dispatch adds constant overhead
that is amortized at longer contexts and on larger models.

\paragraph{4-bit quantized model.}

\begin{table}[h]
\centering
\begin{tabular}{rcccc}
\toprule
Context & Standard (tok/s) & Foveated (tok/s) & Ratio & Memory \\
\midrule
512 & 135 & 96 & 0.71$\times$ & 2.0$\times$ \\
1K & 131 & 97 & 0.74$\times$ & 2.0$\times$ \\
2K & 121 & 98 & 0.81$\times$ & 2.0$\times$ \\
4K & 107 & 67 & 0.63$\times$ & 2.0$\times$ \\
\bottomrule
\end{tabular}
\caption{Qwen2.5-0.5B-Instruct-4bit on M2 8GB, 100 tokens generated
with spike detection enabled. The ratio improves with context as the
bandwidth advantage of reading compressed data grows.}
\end{table}

The primary value proposition is the 2$\times$ memory compression
enabling longer contexts on memory-constrained devices (8GB Mac with
7B models), rather than raw throughput on small models. The blob is
tracked as a proper MLX input array for correct graph pipelining.

Fused kernel correctness: cosine similarity vs dequant-reference =
\textbf{1.000000} on both 7B (D=128) and Llama-1B (D=64) shapes.

\subsection{Quality}

To isolate kernel-level fidelity from model-specific effects, we
measure attention output cosine similarity on synthetic random queries
against real KV pairs from Qwen2.5-0.5B. KV tensors are shaped to
match 7B GQA dimensions (H$_\text{kv}$=8, D=128) to verify that the
kernel handles production-scale tensor shapes. This tests the
quantization and kernel arithmetic, not end-to-end model quality
(which is evaluated via perplexity in Table~\ref{tab:ppl}).

\begin{table}[h]
\centering
\begin{tabular}{rccc}
\toprule
Context & Cosine vs Exact & MAE & Compression \\
\midrule
512 & 0.9956 & 0.0053 & 2.03$\times$ \\
1K & 0.9950 & 0.0040 & 2.02$\times$ \\
4K & 0.9952 & 0.0020 & 2.02$\times$ \\
8K & 0.9954 & 0.0014 & 2.02$\times$ \\
16K & 0.9953 & 0.0010 & 2.02$\times$ \\
32K & 0.9954 & 0.0007 & 2.02$\times$ \\
\bottomrule
\end{tabular}
\caption{Kernel-level foveated attention fidelity at 10\% near tier,
measured on synthetic queries against real model KV pairs shaped to
7B GQA dimensions. Consistent 2.02$\times$ compression across all
context lengths.}
\end{table}

\subsection{Perplexity}

End-to-end perplexity on diverse multi-topic text:

\begin{table}[h]
\centering
\begin{tabular}{llccc}
\toprule
Model & Context & Standard PPL & Foveated PPL & Ratio \\
\midrule
Qwen2.5-0.5B-4bit & 1K & 6.86 & 7.03 & 1.025$\times$ \\
Qwen2.5-0.5B-4bit & 2K & 15.14 & 15.12 & 0.999$\times$ \\
Qwen2.5-0.5B-4bit & 4K & 29.17 & 29.36 & 1.007$\times$ \\
Llama-3.2-1B-4bit & 1K & 8.12 & 8.11 & 0.999$\times$ \\
Llama-3.2-1B-4bit & 2K & 12.45 & 12.46 & 1.001$\times$ \\
\bottomrule
\end{tabular}
\caption{Perplexity within 2.5\% of standard across both model families.
Foveated compression does not measurably degrade language modeling
quality. Note: these are 4-bit weight-quantized models, so absolute PPL
is higher than fp16 baselines; the relevant metric is the
foveated/standard \emph{ratio}, which isolates the effect of KV cache
compression from weight quantization.}
\label{tab:ppl}
\end{table}

\subsection{LongBench-Lite}

We evaluate on 6 representative LongBench~\cite{longbench} tasks (one
per category, 10 samples each) using official THUDM scoring. Context is
truncated to 2K tokens due to memory constraints on the 8GB test
machine.

\begin{table}[h]
\centering
\begin{tabular}{llcc}
\toprule
Task & Category & Standard & Foveated 10/90 \\
\midrule
qasper & Single-doc QA & 2.0 & 1.6 \\
hotpotqa & Multi-doc QA & 0.0 & 0.0 \\
qmsum & Summarization & 5.3 & 7.0 \\
triviaqa & Few-shot & 1.0 & 1.0 \\
passage\_retrieval\_en & Synthetic & 0.0 & 0.0 \\
lcc & Code & 27.7 & 28.1 \\
\textbf{Average} & & \textbf{6.0} & \textbf{6.3} \\
\bottomrule
\end{tabular}
\caption{LongBench-Lite results (Qwen2.5-0.5B-Instruct-4bit, 2K
context). Absolute scores are low due to the small 4-bit model and
context truncation; the relevant comparison is foveated vs standard on
the same model. Foveated matches or slightly exceeds standard on every
task, confirming no quality degradation from KV cache compression.}
\label{tab:longbench}
\end{table}

\subsection{Promotion Recovery}

We evaluate the spike detection and promotion mechanism
(Section~\ref{sec:promotion}) on a multi-fact retrieval benchmark: a
biographical needle containing 4 checkable facts (dates, achievements)
is buried in $\sim$2K tokens of filler text. With 5\% near tier, the
needle lands in the far tier. We compare three modes:

\begin{table}[h]
\centering
\begin{tabular}{lcc}
\toprule
Method & Facts Retrieved & Score \\
\midrule
Standard (fp16 cache) & 56/64 & 88\% \\
Foveated + promotion & 56/64 & 88\% \\
Foveated, no promotion & 32/64 & 50\% \\
\bottomrule
\end{tabular}
\caption{Promotion recovers 24 additional facts (38\% of total),
matching the fp16 baseline. 100\% reproducible across 8 seeds at
depth=0.3. Model: Qwen2.5-0.5B-Instruct-4bit.}
\label{tab:promotion}
\end{table}

Without promotion, far-tier quantization degrades multi-token retrieval
because the model must attend to the same facts across multiple decode
steps. Spike detection identifies these tokens on the first decode step,
promotion restores them to fp16, and subsequent steps benefit from the
higher-quality representation.

This mechanism has no analogue in prior asymmetric quantization
work~\cite{kivi, leankv, kvadaquant, asymkv}, which apply static
quantization without the ability to recover precision at inference time.

\subsection{Key Precision Ablation}

To validate our choice of fp8~E4M3 for keys (following
LeanKV's~\cite{leankv} K8V4 recommendation), we measure single-layer
attention cosine error at 4K context with different key formats (value
format held constant at INT4):

\begin{table}[h]
\centering
\begin{tabular}{lccc}
\toprule
Key Format & Cosine vs Exact & Cosine Error & vs fp8 \\
\midrule
fp16 (exact) & 1.000000 & 0.000000 & --- \\
INT8 per-token & 0.999991 & 0.000009 & $1.0\times$ \\
fp8 E4M3 (ours) & 0.999991 & 0.000009 & $1.0\times$ \\
INT4 per-token & 0.999683 & 0.000317 & $34.3\times$ \\
\bottomrule
\end{tabular}
\caption{Key quantization format comparison at 4K context. fp8~E4M3 and
INT8 are essentially equivalent; INT4 keys produce $34\times$ more
error. We prefer fp8 over INT8 because the 256-entry LUT enables a
single shared-memory read per element vs per-channel scale/zero
arithmetic for INT8.}
\label{tab:key_ablation}
\end{table}

\subsection{Second Model Family: Llama-3.2-1B}

To verify generalization beyond the Qwen architecture, we evaluate on
Llama-3.2-1B-Instruct-4bit (GQA: H$_\text{q}$=32, H$_\text{kv}$=8,
D=64):

\begin{itemize}
\item Fused kernel vs dequant-reference cosine: \textbf{1.000000}
  (kernel perfectly correct on Llama shapes).
\item Foveated vs exact fp16 attention cosine: \textbf{0.994}
  (consistent with Qwen results at equivalent near\%).
\item Short-context token match: 100\% (identical output to standard
  cache at 512 tokens).
\item Throughput: 92.8 tok/s on M2 8GB.
\end{itemize}

\subsection{Tier Ratio Ablation}

We vary the near-tier percentage from 2\% to 100\% and measure
single-layer attention cosine similarity vs exact fp16 at 2K context
(Qwen2.5-0.5B-Instruct-4bit, layer 0):

\begin{table}[ht]
\centering
\begin{tabular}{rrrcc}
\toprule
Near \% & Near Tokens & Far Tokens & Cosine & Error \\
\midrule
2\%  &   40 & 2008 & 0.9766 & 0.0234 \\
5\%  &  102 & 1946 & 0.9844 & 0.0156 \\
10\% &  204 & 1844 & 0.9901 & 0.0099 \\
15\% &  307 & 1741 & 0.9913 & 0.0087 \\
20\% &  409 & 1639 & 0.9938 & 0.0062 \\
30\% &  614 & 1434 & 0.9952 & 0.0048 \\
50\% & 1024 & 1024 & 0.9970 & 0.0030 \\
75\% & 1536 &  512 & 0.9982 & 0.0018 \\
100\% & 2048 &   0 & 1.0000 & 0.0000 \\
\bottomrule
\end{tabular}
\caption{Attention fidelity vs near-tier percentage at 2K context.
The default 10\% achieves 0.990 cosine with 2$\times$ compression.
Quality degrades gracefully: even at 2\% near, cosine remains above
0.976.}
\label{tab:tier_ratio}
\end{table}

The relationship between near percentage and cosine error is
sub-linear: doubling the near tier from 10\% to 20\% only reduces
error by 37\% (0.0099 to 0.0062). This suggests that the most
important tokens are captured even at small near percentages,
with diminishing returns from adding more.

\subsection{Spike Detection Analysis}

We instrument the fused kernel to measure raw spike frequency during
generation (~2K context, 5\% near tier, 80 decode steps).

\paragraph{Frequency.} The kernel detects spikes on 100\% of decode
steps, averaging 182 raw spikes per step across 24 layers $\times$ 8
heads (192 head-layer slots). This corresponds to a 95\% spike rate
per head-layer slot.

\paragraph{Effective rate.} Of 15,748 raw spikes over 80 steps, only
494 (3.1\%) result in actual promotions after the filtering pipeline.
The 32$\times$ reduction from raw to effective demonstrates that
naive promotion would defeat tiered compression.

\paragraph{Layer distribution.} Spikes are distributed roughly
uniformly across layers (4--6\% each), with slightly higher rates
in middle layers (layer 16: 6.1\%) and the final layer (layer 23:
5.8\%). No single layer dominates, confirming that per-layer
independent spike detection is the right design.

\section{Related Work}

\paragraph{KV cache eviction.}
H2O~\cite{h2o} uses cumulative attention scores to identify heavy
hitters and prunes the rest. StreamingLLM~\cite{streamingllm} keeps
attention sinks and a sliding window. SnapKV~\cite{snapkv} selects
important tokens based on observation-window attention patterns. All
three remove tokens from the softmax sum, changing the function being
computed. FoveatedKV retains all tokens.

\paragraph{Asymmetric K/V quantization.}
KIVI~\cite{kivi} first demonstrated that keys and values require
different quantization strategies (per-channel for keys, per-token for
values) and proposed tuning-free INT2/INT4 quantization.
LeanKV~\cite{leankv} made the bit-width asymmetry explicit, proposing
K8V4 configurations and combining them with importance-adaptive
token-level compression and per-head dynamic sparsity.
KV-AdaQuant~\cite{kvadaquant} provided systematic theoretical
justification through singular value and spectral norm analysis,
proposing K4V2 allocations. AsymKV~\cite{asymkv} from Alibaba
conducted systematic examination of attention output error and proposed
distinct configurations for key and value matrices.
KVSplit~\cite{kvsplit} implements K8V4 differentiated precision
specifically for Apple Silicon with Metal support, making it the closest
existing system to FoveatedKV on the same hardware target. KVSplit
applies uniform K8V4 quantization to all tokens; FoveatedKV differs in
three ways: (a)~importance-adaptive tiering (the top 10\% of tokens
retain fp16, concentrating quantization noise on low-attention tokens
rather than applying it uniformly), (b)~a fused kernel that computes
attention across both tiers in a single dispatch with integrated spike
detection, and (c)~the promotion pipeline that dynamically recovers
full-precision representations from disk when the model's attention
pattern shifts. The tiering and promotion are complementary to the
asymmetric precision that KVSplit and FoveatedKV share.

FoveatedKV adopts the established fp8/INT4 asymmetric precision and
contributes: (a)~the foveated two-tier architecture with per-head
importance-adaptive partitioning, (b)~the fused Metal kernel with
integrated spike detection at zero overhead, and (c)~the promotion
pipeline that dynamically recovers fp16 precision from NVMe-backed
archives---a capability not present in any of the above systems.

\paragraph{Foveated rendering.}
Foveated rendering in VR~\cite{guenter2012} renders the gaze center at
full resolution and reduces quality in the periphery. We adapt this
principle to attention: the ``gaze'' is the query vector, the ``fovea''
is the near tier, and the ``periphery'' is the far tier.

\section{Limitations}

\paragraph{Model scale.} Quality and throughput results are on
Qwen2.5-0.5B and Llama-3.2-1B; kernel microbenchmarks use 7B-equivalent
tensor shapes. Both Qwen (GQA 8-head, D=128) and Llama (GQA 8-head,
D=64) architectures show consistent behavior. We have not evaluated on
13B+ models. Prior work on asymmetric K/V quantization~\cite{leankv,
kvadaquant} has validated the approach at larger scales; our system-level
contributions (kernel, promotion) need separate validation.

\paragraph{Evaluation breadth.} End-to-end quality is evaluated on
LongBench-Lite (Table~\ref{tab:longbench}), perplexity, and a custom
multi-fact retrieval benchmark. LongBench results are on a 0.5B model
with 2K context truncation; full-length evaluation on larger models
with RULER or complete LongBench would strengthen the claims.

\paragraph{Hardware.} All results are on Apple M2 8GB. End-to-end
throughput is measured on a 0.5B model where inference is not
swap-bound; the speedups (1.03--1.45$\times$) reflect genuine bandwidth
savings rather than memory-pressure effects. On machines with higher
memory bandwidth, the absolute tok/s would increase but the relative
speedup should persist since it derives from reading 2$\times$ fewer KV
bytes. The kernel bandwidth advantage (up to 3.34$\times$ at 16K on 7B
shapes) is hardware-independent.

\paragraph{Comparison scope.} We do not directly benchmark against
LeanKV~\cite{leankv} or KV-AdaQuant~\cite{kvadaquant}. Kernel-level
comparison is infeasible (those systems target CUDA), but the
algorithmic components---tier assignment policy, compression quality at
equivalent bit-budgets---could be compared via Python reference
implementations on the same models and benchmarks. A controlled
algorithmic comparison, and a head-to-head with llama.cpp's quantized
KV cache on the same Mac hardware, would meaningfully strengthen the
evaluation. We leave these for future work.

\section{Conclusion}

FoveatedKV demonstrates that combining established asymmetric K/V
quantization with importance-adaptive tiering, a fused Metal kernel, and
spike-driven promotion yields a practical system for Apple Silicon
inference that is \emph{faster} than standard fp16 at all context
lengths. On an M2 8GB Mac, end-to-end throughput improves from
1.03$\times$ at 512 tokens to 1.45$\times$ at 4K context, because the
fused kernel reads 2$\times$ fewer bytes from memory. At the kernel
level, speedups reach 3.34$\times$ at 16K on 7B-equivalent shapes. The
system achieves 2$\times$ memory compression with negligible quality loss
(0.995+ cosine fidelity, non-accumulating perplexity ratio). Direct
attention module patching eliminates Python overhead, and tracking the
compressed blob as an MLX input array enables proper graph pipelining.
The promotion mechanism---detecting spikes during online softmax at zero
overhead and recovering exact fp16 from NVMe-backed archives---closes the
quality gap for sustained multi-fact retrieval, a capability not present
in prior quantization-only approaches.

\begin{thebibliography}{99}
\bibitem{h2o} Zhang et al., ``H2O: Heavy-Hitter Oracle for Efficient Generative Inference of Large Language Models,'' NeurIPS 2023.
\bibitem{snapkv} Li et al., ``SnapKV: LLM Knows What You Are Looking For Before Generation,'' 2024.
\bibitem{streamingllm} Xiao et al., ``Efficient Streaming Language Models with Attention Sinks,'' ICLR 2024.
\bibitem{kivi} Liu et al., ``KIVI: A Tuning-Free Asymmetric 2bit Quantization for KV Cache,'' ICML 2024.
\bibitem{leankv} Chen et al., ``LeanKV: KV Cache Compression with Differentiated Key-Value Precision,'' arXiv:2412.03131, 2024.
\bibitem{kvadaquant} Li et al., ``KV-AdaQuant: Adaptive KV Cache Quantization with Key-Value Precision Allocation,'' arXiv:2502.15075, 2025.
\bibitem{asymkv} Li et al., ``AsymKV: Asymmetric Key-Value Cache Quantization,'' arXiv:2410.13212, 2024.
\bibitem{kvsplit} Paul, ``KVSplit: Differentiated Key-Value Precision for Apple Silicon,'' \url{https://github.com/dipampaul17/KVSplit}, 2025.
\bibitem{longbench} Bai et al., ``LongBench: A Bilingual, Multitask Benchmark for Long Context Understanding,'' ACL 2024.
\bibitem{guenter2012} Guenter et al., ``Foveated 3D Graphics,'' ACM SIGGRAPH Asia 2012.
\end{thebibliography}

\end{document}
